% --- Archon physics formalization source begin ---
% archon:physics
% archon:covers PhyXMiniProblems/problem_phyx_mini_0134.lean
% archon:source-report reports/phyx_mini/problem_phyx_mini_0134.source.json

\chapter{Physics problem phyx\_mini\_0134}
\label{ch:PhyXMiniProblems_problem_phyx_mini_0134}

\paragraph{Problem source.}
\#\# Physical scenario

Two converging lenses, A and B, with focal lengths \textbackslash{}(f\_A = 20.0\textbackslash{},\textbackslash{}mathrm\{cm\}\textbackslash{}) and \textbackslash{}(f\_B = 25.0\textbackslash{},\textbackslash{}mathrm\{cm\}\textbackslash{}), are placed \textbackslash{}(80.0\textbackslash{},\textbackslash{}mathrm\{cm\}\textbackslash{}) apart, as shown in figure. An object is placed \textbackslash{}(60.0\textbackslash{},\textbackslash{}mathrm\{cm\}\textbackslash{}) in front of the first lens as shown in figure.

\#\# Question

Determine the magnification of the final image formed by the combination of the two lenses.

\#\# Answer choices

- A: A: -0.30

- B: B: -0.40

- C: C: -0.50

- D: D: -0.60

\#\# Figure caption (auxiliary; use the image as primary evidence)

The image consists of two main sections illustrating a lens setup with ray diagrams.

**Top Diagram:**
- **Lenses:** There are two converging lenses, labeled as Lens A and Lens B.
- **Focal Points:** Each lens has focal points labeled \textbackslash{}( F'\_A \textbackslash{}), \textbackslash{}( F\_A \textbackslash{}), \textbackslash{}( F'\_B \textbackslash{}), and \textbackslash{}( F\_B \textbackslash{}). 
- **Distance:** The distance between Lens A and Lens B is marked as \textbackslash{}( 80.0 \textbackslash{}, \textbackslash{}text\{cm\} \textbackslash{}).

**Bottom Diagram:**
- **Object:** An upright object, represented by an arrow, is placed to the left of Lens A at point \textbackslash{}( O \textbackslash{}).
- **Ray Diagrams:** Multiple rays are traced from the object through both lenses.
  - Ray 1 is parallel to the principal axis and refracts through \textbackslash{}( F\_A \textbackslash{}).
  - Ray 2 passes through the optical center of Lens A.
  - Ray 3 passes through \textbackslash{}( F'\_A \textbackslash{}) and refracts parallel to the principal axis.
  - Ray 4 travels directly along the principal axis.
  - Similar rays (labeled 1', 2', 3') are shown interacting with Lens B.
- **Images:** 
  - Image \textbackslash{}( I\_A \textbackslash{}) is formed by Lens A before Lens B.
  - \textbackslash{}( O\_B \textbackslash{}) is depicted at the same location as \textbackslash{}( I\_A \textbackslash{}), labeled as \textbackslash{}( (= I\_A) \textbackslash{}).
  - Final image \textbackslash{}( I\_B \textbackslash{}) is formed to the right of Lens B.
- **Distances:**
  - \textbackslash{}( d\_\{oA\} \textbackslash{}) and \textbackslash{}( d\_\{iA\} \textbackslash{}) represent object and image distances for Lens A.
  - \textbackslash{}( d\_\{oB\} \textbackslash{}) and \textbackslash{}( d\_\{iB\} \textbackslash{}) are the object and image distances for Lens B.

Overall, the diagram illustrates the path of light through a two-lens system and shows how the final image is formed in relation to both lenses.

\#\# Dataset metadata

Geometrical Optics; Multi-Formula Reasoning, Spatial Relation Reasoning

\paragraph{Recorded answer/context.}
C: C: -0.50

\paragraph{Figure/image path.}
/root/proposal\_for\_physic/science-mango/phyx\_mini\_run/phyx\_data/test\_image/134.png

\paragraph{Formalization target.}
create a compiling Lean file with sorry bodies at `PhyXMiniProblems/problem\_phyx\_mini\_0134.lean`. The Lean declarations must preserve the physical quantities, dimensions or dimensional roles, figure labels, governing-law hypotheses, and final relation expressed by this problem.
Use Mathlib/Physlib names found through LeanExplore where available. If a physics API is missing, introduce faithful local abstractions rather than scalar placeholder aliases.

\begin{theorem}[Physics formalization target]
\label{thm:physics:phyx_mini_0134:target}
\lean{PhyXMiniProblems.ProblemPhyXMini0134.problem_phyx_mini_0134}
\uses{def:physics:phyx-mini-0134:phyxminiproblems-problemphyxmini0134-twolensdiagram, def:physics:phyx-mini-0134:phyxminiproblems-problemphyxmini0134-matchesproblemdata, def:physics:phyx-mini-0134:phyxminiproblems-problemphyxmini0134-matchestwolensfigure, def:physics:phyx-mini-0134:phyxminiproblems-problemphyxmini0134-hasphysicaltwolensconfiguration, def:physics:phyx-mini-0134:phyxminiproblems-problemphyxmini0134-satisfiestwolensparaxiallaws}
The assigned autoformalize agent should translate this physics problem into Lean declarations in the covered file, with theorem and lemma proof bodies written as `by sorry`.
\end{theorem}
\begin{proof}
This is an autoformalization task, not a proof task. Produce statements that can later be proved by the physics prover without weakening the physical model.
\end{proof}

% --- Archon declaration topology begin ---
\section{Declaration topology}
\begin{definition}[Optical Length]
  \label{def:physics:phyx-mini-0134:phyxminiproblems-problemphyxmini0134-opticallength}
  \lean{PhyXMiniProblems.ProblemPhyXMini0134.OpticalLength}
  A signed physical optical length, independent of the unit used to read it
\end{definition}
\begin{proof}
  This is the declared model definition of Optical Length.
\end{proof}
\begin{definition}[unit Choices For Length]
  \label{def:physics:phyx-mini-0134:phyxminiproblems-problemphyxmini0134-unitchoicesforlength}
  \lean{PhyXMiniProblems.ProblemPhyXMini0134.unitChoicesForLength}
  Unit choices obtained from SI by replacing only the length unit
\end{definition}
\begin{proof}
  This is the declared model definition of unit Choices For Length.
\end{proof}
\begin{definition}[length Readout]
  \label{def:physics:phyx-mini-0134:phyxminiproblems-problemphyxmini0134-lengthreadout}
  \lean{PhyXMiniProblems.ProblemPhyXMini0134.lengthReadout}
  \uses{def:physics:phyx-mini-0134:phyxminiproblems-problemphyxmini0134-opticallength, def:physics:phyx-mini-0134:phyxminiproblems-problemphyxmini0134-unitchoicesforlength}
  Scalar readout of a physical optical length in a selected length unit
\end{definition}
\begin{proof}
  This is the declared model definition of length Readout.
\end{proof}
\begin{definition}[Lens Label]
  \label{def:physics:phyx-mini-0134:phyxminiproblems-problemphyxmini0134-lenslabel}
  \lean{PhyXMiniProblems.ProblemPhyXMini0134.LensLabel}
  The two lenses in the order traversed by the light
\end{definition}
\begin{proof}
  This is the declared model definition of Lens Label.
\end{proof}
\begin{definition}[Thin Lens Kind]
  \label{def:physics:phyx-mini-0134:phyxminiproblems-problemphyxmini0134-thinlenskind}
  \lean{PhyXMiniProblems.ProblemPhyXMini0134.ThinLensKind}
  Qualitative focal-length sign role of a thin lens
\end{definition}
\begin{proof}
  This is the declared model definition of Thin Lens Kind.
\end{proof}
\begin{definition}[Optical Approximation]
  \label{def:physics:phyx-mini-0134:phyxminiproblems-problemphyxmini0134-opticalapproximation}
  \lean{PhyXMiniProblems.ProblemPhyXMini0134.OpticalApproximation}
  Optical approximation used by the ray diagram
\end{definition}
\begin{proof}
  This is the declared model definition of Optical Approximation.
\end{proof}
\begin{definition}[Axis Point]
  \label{def:physics:phyx-mini-0134:phyxminiproblems-problemphyxmini0134-axispoint}
  \lean{PhyXMiniProblems.ProblemPhyXMini0134.AxisPoint}
  Labeled points on the principal axis in the primary figure
\end{definition}
\begin{proof}
  This is the declared model definition of Axis Point.
\end{proof}
\begin{definition}[Arrow Label]
  \label{def:physics:phyx-mini-0134:phyxminiproblems-problemphyxmini0134-arrowlabel}
  \lean{PhyXMiniProblems.ProblemPhyXMini0134.ArrowLabel}
  Object and image arrows whose transverse orientations are shown
\end{definition}
\begin{proof}
  This is the declared model definition of Arrow Label.
\end{proof}
\begin{definition}[Image Orientation]
  \label{def:physics:phyx-mini-0134:phyxminiproblems-problemphyxmini0134-imageorientation}
  \lean{PhyXMiniProblems.ProblemPhyXMini0134.ImageOrientation}
  Transverse orientation relative to the upright source arrow
\end{definition}
\begin{proof}
  This is the declared model definition of Image Orientation.
\end{proof}
\begin{definition}[Image Nature]
  \label{def:physics:phyx-mini-0134:phyxminiproblems-problemphyxmini0134-imagenature}
  \lean{PhyXMiniProblems.ProblemPhyXMini0134.ImageNature}
  Whether the outgoing rays themselves meet at the image point
\end{definition}
\begin{proof}
  This is the declared model definition of Image Nature.
\end{proof}
\begin{definition}[Ray Label]
  \label{def:physics:phyx-mini-0134:phyxminiproblems-problemphyxmini0134-raylabel}
  \lean{PhyXMiniProblems.ProblemPhyXMini0134.RayLabel}
  Principal-ray labels printed in the lower diagram
\end{definition}
\begin{proof}
  This is the declared model definition of Ray Label.
\end{proof}
\begin{definition}[Principal Ray Role]
  \label{def:physics:phyx-mini-0134:phyxminiproblems-problemphyxmini0134-principalrayrole}
  \lean{PhyXMiniProblems.ProblemPhyXMini0134.PrincipalRayRole}
  Qualitative construction rule represented by a principal ray
\end{definition}
\begin{proof}
  This is the declared model definition of Principal Ray Role.
\end{proof}
\begin{definition}[Thin Lens]
  \label{def:physics:phyx-mini-0134:phyxminiproblems-problemphyxmini0134-thinlens}
  \lean{PhyXMiniProblems.ProblemPhyXMini0134.ThinLens}
  \uses{def:physics:phyx-mini-0134:phyxminiproblems-problemphyxmini0134-opticallength, def:physics:phyx-mini-0134:phyxminiproblems-problemphyxmini0134-thinlenskind}
  A thin lens carrying its physical kind and signed focal length
\end{definition}
\begin{proof}
  This is the declared model definition of Thin Lens.
\end{proof}
\begin{definition}[Two Lens Diagram]
  \label{def:physics:phyx-mini-0134:phyxminiproblems-problemphyxmini0134-twolensdiagram}
  \lean{PhyXMiniProblems.ProblemPhyXMini0134.TwoLensDiagram}
  \uses{def:physics:phyx-mini-0134:phyxminiproblems-problemphyxmini0134-opticallength, def:physics:phyx-mini-0134:phyxminiproblems-problemphyxmini0134-lenslabel, def:physics:phyx-mini-0134:phyxminiproblems-problemphyxmini0134-opticalapproximation, def:physics:phyx-mini-0134:phyxminiproblems-problemphyxmini0134-axispoint, def:physics:phyx-mini-0134:phyxminiproblems-problemphyxmini0134-arrowlabel, def:physics:phyx-mini-0134:phyxminiproblems-problemphyxmini0134-imageorientation, def:physics:phyx-mini-0134:phyxminiproblems-problemphyxmini0134-imagenature, def:physics:phyx-mini-0134:phyxminiproblems-problemphyxmini0134-raylabel, def:physics:phyx-mini-0134:phyxminiproblems-problemphyxmini0134-principalrayrole, def:physics:phyx-mini-0134:phyxminiproblems-problemphyxmini0134-thinlens}
  Physical quantities and labels of the two-lens construction. signedImageDistance uses the Cartesian convention: it is positive for the two real images drawn to the right of their respective lenses. The stage and net transverse magnifications retain their signs and are dimensionless
\end{definition}
\begin{proof}
  This is the declared model definition of Two Lens Diagram.
\end{proof}
\begin{definition}[focal Length Readout]
  \label{def:physics:phyx-mini-0134:phyxminiproblems-problemphyxmini0134-focallengthreadout}
  \lean{PhyXMiniProblems.ProblemPhyXMini0134.focalLengthReadout}
  \uses{def:physics:phyx-mini-0134:phyxminiproblems-problemphyxmini0134-lengthreadout, def:physics:phyx-mini-0134:phyxminiproblems-problemphyxmini0134-lenslabel, def:physics:phyx-mini-0134:phyxminiproblems-problemphyxmini0134-twolensdiagram}
  Focal-length readout for one lens
\end{definition}
\begin{proof}
  This is the declared model definition of focal Length Readout.
\end{proof}
\begin{definition}[object Distance Readout]
  \label{def:physics:phyx-mini-0134:phyxminiproblems-problemphyxmini0134-objectdistancereadout}
  \lean{PhyXMiniProblems.ProblemPhyXMini0134.objectDistanceReadout}
  \uses{def:physics:phyx-mini-0134:phyxminiproblems-problemphyxmini0134-lengthreadout, def:physics:phyx-mini-0134:phyxminiproblems-problemphyxmini0134-lenslabel, def:physics:phyx-mini-0134:phyxminiproblems-problemphyxmini0134-twolensdiagram}
  Object-distance readout for one lens
\end{definition}
\begin{proof}
  This is the declared model definition of object Distance Readout.
\end{proof}
\begin{definition}[image Distance Readout]
  \label{def:physics:phyx-mini-0134:phyxminiproblems-problemphyxmini0134-imagedistancereadout}
  \lean{PhyXMiniProblems.ProblemPhyXMini0134.imageDistanceReadout}
  \uses{def:physics:phyx-mini-0134:phyxminiproblems-problemphyxmini0134-lengthreadout, def:physics:phyx-mini-0134:phyxminiproblems-problemphyxmini0134-lenslabel, def:physics:phyx-mini-0134:phyxminiproblems-problemphyxmini0134-twolensdiagram}
  Signed image-distance readout for one lens
\end{definition}
\begin{proof}
  This is the declared model definition of image Distance Readout.
\end{proof}
\begin{definition}[axis Position Readout]
  \label{def:physics:phyx-mini-0134:phyxminiproblems-problemphyxmini0134-axispositionreadout}
  \lean{PhyXMiniProblems.ProblemPhyXMini0134.axisPositionReadout}
  \uses{def:physics:phyx-mini-0134:phyxminiproblems-problemphyxmini0134-lengthreadout, def:physics:phyx-mini-0134:phyxminiproblems-problemphyxmini0134-axispoint, def:physics:phyx-mini-0134:phyxminiproblems-problemphyxmini0134-twolensdiagram}
  Coordinate readout of a labeled point on the principal axis
\end{definition}
\begin{proof}
  This is the declared model definition of axis Position Readout.
\end{proof}
\begin{definition}[axial Separation Readout]
  \label{def:physics:phyx-mini-0134:phyxminiproblems-problemphyxmini0134-axialseparationreadout}
  \lean{PhyXMiniProblems.ProblemPhyXMini0134.axialSeparationReadout}
  \uses{def:physics:phyx-mini-0134:phyxminiproblems-problemphyxmini0134-axispoint, def:physics:phyx-mini-0134:phyxminiproblems-problemphyxmini0134-twolensdiagram, def:physics:phyx-mini-0134:phyxminiproblems-problemphyxmini0134-axispositionreadout}
  Directed axial separation from start to finish
\end{definition}
\begin{proof}
  This is the declared model definition of axial Separation Readout.
\end{proof}
\begin{definition}[lens Separation Readout]
  \label{def:physics:phyx-mini-0134:phyxminiproblems-problemphyxmini0134-lensseparationreadout}
  \lean{PhyXMiniProblems.ProblemPhyXMini0134.lensSeparationReadout}
  \uses{def:physics:phyx-mini-0134:phyxminiproblems-problemphyxmini0134-twolensdiagram, def:physics:phyx-mini-0134:phyxminiproblems-problemphyxmini0134-axialseparationreadout}
  Separation from lens A to lens B
\end{definition}
\begin{proof}
  This is the declared model definition of lens Separation Readout.
\end{proof}
\begin{definition}[Matches Problem Data]
  \label{def:physics:phyx-mini-0134:phyxminiproblems-problemphyxmini0134-matchesproblemdata}
  \lean{PhyXMiniProblems.ProblemPhyXMini0134.MatchesProblemData}
  \uses{def:physics:phyx-mini-0134:phyxminiproblems-problemphyxmini0134-twolensdiagram, def:physics:phyx-mini-0134:phyxminiproblems-problemphyxmini0134-focallengthreadout, def:physics:phyx-mini-0134:phyxminiproblems-problemphyxmini0134-objectdistancereadout, def:physics:phyx-mini-0134:phyxminiproblems-problemphyxmini0134-lensseparationreadout}
  Numerical data stated in the problem: f\_A = 20 cm, f\_B = 25 cm, lens separation 80 cm, and original object distance d\_oA = 60 cm
\end{definition}
\begin{proof}
  This is the declared model definition of Matches Problem Data.
\end{proof}
\begin{definition}[Matches Two Lens Figure]
  \label{def:physics:phyx-mini-0134:phyxminiproblems-problemphyxmini0134-matchestwolensfigure}
  \lean{PhyXMiniProblems.ProblemPhyXMini0134.MatchesTwoLensFigure}
  \uses{def:physics:phyx-mini-0134:phyxminiproblems-problemphyxmini0134-twolensdiagram, def:physics:phyx-mini-0134:phyxminiproblems-problemphyxmini0134-focallengthreadout, def:physics:phyx-mini-0134:phyxminiproblems-problemphyxmini0134-objectdistancereadout, def:physics:phyx-mini-0134:phyxminiproblems-problemphyxmini0134-imagedistancereadout, def:physics:phyx-mini-0134:phyxminiproblems-problemphyxmini0134-axispositionreadout, def:physics:phyx-mini-0134:phyxminiproblems-problemphyxmini0134-axialseparationreadout}
  Literal qualitative labels and axial identifications read from the figure. No numerical image distance or magnification is asserted here
\end{definition}
\begin{proof}
  This is the declared model definition of Matches Two Lens Figure.
\end{proof}
\begin{definition}[Has Physical Two Lens Configuration]
  \label{def:physics:phyx-mini-0134:phyxminiproblems-problemphyxmini0134-hasphysicaltwolensconfiguration}
  \lean{PhyXMiniProblems.ProblemPhyXMini0134.HasPhysicalTwoLensConfiguration}
  \uses{def:physics:phyx-mini-0134:phyxminiproblems-problemphyxmini0134-lenslabel, def:physics:phyx-mini-0134:phyxminiproblems-problemphyxmini0134-twolensdiagram, def:physics:phyx-mini-0134:phyxminiproblems-problemphyxmini0134-focallengthreadout, def:physics:phyx-mini-0134:phyxminiproblems-problemphyxmini0134-objectdistancereadout, def:physics:phyx-mini-0134:phyxminiproblems-problemphyxmini0134-imagedistancereadout, def:physics:phyx-mini-0134:phyxminiproblems-problemphyxmini0134-axispositionreadout}
  Positivity and left-to-right ordering for the physical branch shown
\end{definition}
\begin{proof}
  This is the declared model definition of Has Physical Two Lens Configuration.
\end{proof}
\begin{definition}[Satisfies Two Lens Paraxial Laws]
  \label{def:physics:phyx-mini-0134:phyxminiproblems-problemphyxmini0134-satisfiestwolensparaxiallaws}
  \lean{PhyXMiniProblems.ProblemPhyXMini0134.SatisfiesTwoLensParaxialLaws}
  \uses{def:physics:phyx-mini-0134:phyxminiproblems-problemphyxmini0134-lenslabel, def:physics:phyx-mini-0134:phyxminiproblems-problemphyxmini0134-twolensdiagram, def:physics:phyx-mini-0134:phyxminiproblems-problemphyxmini0134-focallengthreadout, def:physics:phyx-mini-0134:phyxminiproblems-problemphyxmini0134-objectdistancereadout, def:physics:phyx-mini-0134:phyxminiproblems-problemphyxmini0134-imagedistancereadout}
  Standard thin-lens, signed transverse-magnification, and composition laws. The Gaussian equation is written without division as f (d\_o + d\_i) = d\_o d\_i. The stage law is m d\_o = -d\_i, and the net magnification is the product of the two stage magnifications. All dimensionful factors in an equation are read in the same arbitrary length unit
\end{definition}
\begin{proof}
  This is the declared model definition of Satisfies Two Lens Paraxial Laws.
\end{proof}
\begin{lemma}[derived stage distances and magnifications]
  \label{lem:physics:phyx-mini-0134:phyxminiproblems-problemphyxmini0134-derived-stage-distances-and-magnifications}
  \lean{PhyXMiniProblems.ProblemPhyXMini0134.derived_stage_distances_and_magnifications}
  \uses{def:physics:phyx-mini-0134:phyxminiproblems-problemphyxmini0134-twolensdiagram, def:physics:phyx-mini-0134:phyxminiproblems-problemphyxmini0134-objectdistancereadout, def:physics:phyx-mini-0134:phyxminiproblems-problemphyxmini0134-imagedistancereadout, def:physics:phyx-mini-0134:phyxminiproblems-problemphyxmini0134-matchesproblemdata, def:physics:phyx-mini-0134:phyxminiproblems-problemphyxmini0134-matchestwolensfigure, def:physics:phyx-mini-0134:phyxminiproblems-problemphyxmini0134-hasphysicaltwolensconfiguration, def:physics:phyx-mini-0134:phyxminiproblems-problemphyxmini0134-satisfiestwolensparaxiallaws}
  The two imaging stages determine the unprinted intermediate quantities: d\_iA = 30 cm, d\_oB = 50 cm, d\_iB = 50 cm, m\_A = -1/2, and m\_B = -1
\end{lemma}
\begin{proof}
  The conclusion follows from the governing relations and figure readouts named in the statement of derived stage distances and magnifications.
\end{proof}
\begin{definition}[Answer Choice]
  \label{def:physics:phyx-mini-0134:phyxminiproblems-problemphyxmini0134-answerchoice}
  \lean{PhyXMiniProblems.ProblemPhyXMini0134.AnswerChoice}
  Labels of the four displayed answers
\end{definition}
\begin{proof}
  This is the declared model definition of Answer Choice.
\end{proof}
\begin{definition}[magnification Readout]
  \label{def:physics:phyx-mini-0134:phyxminiproblems-problemphyxmini0134-answerchoice-magnificationreadout}
  \lean{PhyXMiniProblems.ProblemPhyXMini0134.AnswerChoice.magnificationReadout}
  \uses{def:physics:phyx-mini-0134:phyxminiproblems-problemphyxmini0134-answerchoice}
  Signed dimensionless magnification printed beside each answer
\end{definition}
\begin{proof}
  This is the declared model definition of magnification Readout.
\end{proof}
\begin{definition}[recorded Answer Choice]
  \label{def:physics:phyx-mini-0134:phyxminiproblems-problemphyxmini0134-recordedanswerchoice}
  \lean{PhyXMiniProblems.ProblemPhyXMini0134.recordedAnswerChoice}
  \uses{def:physics:phyx-mini-0134:phyxminiproblems-problemphyxmini0134-answerchoice}
  Dataset metadata only: the source records choice C
\end{definition}
\begin{proof}
  This is the declared model definition of recorded Answer Choice.
\end{proof}
\begin{definition}[Matches Displayed Answer]
  \label{def:physics:phyx-mini-0134:phyxminiproblems-problemphyxmini0134-matchesdisplayedanswer}
  \lean{PhyXMiniProblems.ProblemPhyXMini0134.MatchesDisplayedAnswer}
  \uses{def:physics:phyx-mini-0134:phyxminiproblems-problemphyxmini0134-twolensdiagram, def:physics:phyx-mini-0134:phyxminiproblems-problemphyxmini0134-answerchoice}
  Exact agreement between the physical net magnification and a displayed choice
\end{definition}
\begin{proof}
  This is the declared model definition of Matches Displayed Answer.
\end{proof}
% --- Archon declaration topology end ---
% --- Archon physics formalization source end ---
